\section{Human-in-the-Loop and Human Oversight in AI-Assisted Development}

As we already saw, splitting a task across several agents does not remove the risk of an early mistake propagating undetected through the rest of a pipeline --- it only changes where in the process an error is likely to be introduced. This raises a question that applies just as much to a single agent working alone as it does to a coordinated multi-agent pipeline: at what point, and in what way, does a human actually enter this process?

\textbf{Human-in-the-loop as a design choice, not an incidental fact.} It is tempting to describe human oversight loosely, as "a person eventually looked at the result." But in the literature on human-AI collaboration, human-in-the-loop (HITL) refers to something more specific: a deliberate design decision about where, in an otherwise automated process, human judgment is required before the process is allowed to continue, and what form that judgment takes. Human-in-the-loop AI combines automated capability with human oversight, feedback, and decision-making at various stages of a system's operation, rather than treating human involvement as something bolted on only at the very end. This means the relevant question for any AI-assisted workflow is not simply "is a human involved?" but "at which stage, and with what authority?"

\textbf{Forms of oversight.} Human involvement can take noticeably different forms, and these are worth distinguishing rather than lumping together as "review." At one end, an agent's output can be accepted automatically, with no human check at all. At the other end, a process can require explicit human approval before it is allowed to proceed to the next stage --- the process pauses at a defined point and genuinely waits, rather than continuing regardless. In between sits review after the fact: a human examines a finished output once it already exists, and can accept, correct, or discard it, but the work leading up to that point was not gated on their approval. The literature draws a related distinction between keeping a human directly in a decision cycle versus having a human monitor an otherwise autonomous process and step in only when something looks wrong. Applied to the multi-agent example already introduced --- a planner breaking a task into subtasks, an executor implementing one, a critic checking it --- a human could sit at any of several points: approving the plan before any code is written, reviewing each subtask only after the critic has already accepted it, or not reviewing anything until the entire feature is finished.

\textbf{Why placement matters, not just presence.} The distinction between these forms is not merely procedural. If a human's only review happens once every subtask is finished, an error introduced at the very first step --- a plan built on a misunderstood requirement --- has already shaped every subtask built on top of it by the time anyone catches it, and correcting it may mean discarding work that depended on the flawed plan. Reviewing at an earlier gate, such as approving the plan itself before any implementation begins, catches the same class of error while it is still cheap to correct. This mirrors a long-standing observation in traditional software engineering, where a mistake caught during requirements or design is generally far less costly to fix than the same mistake caught only after it has been implemented and built upon. The same logic carries over directly to AI-assisted development: where a human is placed in an agent's workflow determines not just whether an error can be caught, but how much of the surrounding work has already been built on top of it by the time it is.

\textbf{What human judgment contributes that the agent cannot supply on its own.} As already discussed, an agent has no way of verifying its output against anything beyond what was explicitly given to it --- it can check that its code runs, or that it satisfies a stated instruction, but it has no independent way of judging whether that instruction actually captured what was intended. A human reviewer can catch a different class of problem: cases where an agent followed its instructions correctly, and its output is internally consistent and free of obvious errors, but still does not achieve what was actually meant, because the instruction itself was incomplete, ambiguous, or simply did not anticipate a particular situation. This is a failure mode no amount of agent-side verification can catch on its own, since the agent has nothing else to check its work against.

Taken together, human oversight in AI-assisted development is best understood as a variable that can be tuned --- in its placement, its frequency, and the amount of authority it carries at each point --- rather than as a single, fixed ingredient that is either present or absent.

